%%%%%%%%%%  scan idx 338  |  folio 325  %%%%%%%%%%
de \(\underline{E}\)~» ;
pour définir la \(n\)-ième classe de Chern, nous procèderons de façon analogue, en ramenant \(\underline{E}\) au rang \(n\) en considérant \((\cl(\underline{E})-d)+n\), et nous prendrons la somme alternée des puissances extérieures de l'élément ainsi obtenu ; on «~interprêtera~» la proposition qui suit en y faisant \(x=\cl(\underline{E})-d\) :

\textbf{Proposition 3.2.} Soit \(K\) un \(\lambda\)-anneau.
\begin{enumerate}[label=\roman*)]
\item \(\forall x\in K,\ \forall n\in\mathbb N,\quad \lambda^n(x+n-1)=(-1)^n\sum_{i=0}^n(-1)^i\lambda^i(x+n).\)
\item Si on pose \(\gamma^n(x)=\lambda^n(x+n-1)\), et \(\gamma_t(x)=\sum_{n\ge0}\gamma^n(x)t^n\), alors :
\[
  \gamma_t(x)=\lambda_{t/1-t}(x),\qquad\text{donc}
\]
\[
  \lambda_t(x)=\gamma_{t/1+t}(x).
\]
\end{enumerate}

On voit donc que la connaissance des opérations \(\gamma\) équivaut à celle des opérations \(\lambda\).

\emph{Démonstration.}

i) On a
\[
  \lambda^n(x+n-1)=\sum_{i=0}^{n}\lambda^i(x+n)\lambda^{n-i}(-1).
\]
De \(\lambda_t(1)=1+t\), on tire : \(\lambda_t(-1)=1/1+t\), d'où : \(\lambda^k(-1)=(-1)^k\), ce qui donne la relation cherchée.

ii) \(\lambda_{t/1-t}(x)=\sum_{n\ge0}\lambda^n(x).t^n(1-t)^{-n}.\)
On développe \(t^n(1-t)^{-n}=t^n\sum_{k\ge0}\binom{-n}{k}(-1)^kt^k=\sum_{p\ge n}\binom{p-1}{p-n}t^p\). \(\lambda_{t/1-t}(x)=\sum_{n\ge0}\bigl[\lambda^n(x).\sum_{p\ge n}\binom{p-1}{p-n}t^p\bigr]\)~; le coefficient de \(t^p\) est donc :
\[
  \sum_{n=0}\lambda^n(x).\binom{p-1}{p-n}
  =\sum_{n=0}\lambda^n(x)\lambda^{p-n}(p-1)
  =\lambda^p(x+p-1).
\]

